\section{Why Three Dimensions}

Our space has three dimensions: length, width, height. Most theories simply assume this. They write ``three'' at the top of the page and move on. Vibe Theory does something more satisfying. It does not get to choose. The number three is forced, twice over. This chapter shows how.

\subsection{First squeeze: only 2, 3, and 4 are even possible}

Remember that a crystal is a regular tiling, generated by a kaleidoscope, with finite cells. Ask: in how many dimensions can such a thing exist in hyperbolic space?

In the hyperbolic plane (two dimensions), there are infinitely many, one for every pair of integers that passes the curvature test. Lots of room, lots of crystals.

In hyperbolic three-dimensional space, the number collapses. There are exactly four such crystals. Not infinitely many, four. The dodecagrid \{5, 3, 4\} is one of them.

In four-dimensional hyperbolic space, there are exactly five.

In five dimensions and above, there are none. Zero. It becomes impossible to fill hyperbolic space with finite regular cells at all. The geometry simply will not allow it.

This is not an opinion or a tuning. It is a hard mathematical classification, known independently of this theory, and the model's simulation reproduces it exactly: infinitely many, then four, then five, then nothing, forever. So a universe built as a finite-celled crystal can only have two, three, or four dimensions of space. The door is open just a crack.

\subsection{Second squeeze: only three holds matter together}

That leaves three candidates: two, three, or four. The second argument picks three.

Think about orbits, a planet going around a star, an electron around a nucleus. For stable matter to exist, you need orbits that close up and repeat, so things can stay bound and structured over time. Whether orbits close depends on the number of dimensions, through how gravity and other forces fall off with distance.

\begin{itemize}
\item In two dimensions, an orbit does not quite close. It slowly rotates, drifting around forever, never repeating. No stable, repeating structure.
\item In three dimensions, orbits close perfectly. A planet traces the same ellipse again and again. This is the world of stable atoms and solar systems.
\item In four dimensions and above, orbits are unstable. The slightest nudge and the planet either spirals into the star or flies off to infinity. Nothing stays bound.
\end{itemize}

So of the three dimensions the crystal could possibly have, only three supports the stable, repeating structures that matter and life require. Two is too loose, four is too fragile, three is just right. This is an old argument, going back a century, and the model recovers it on the discrete crystal.

\subsection{Putting it together}

The number of spatial dimensions is not a free parameter. The geometry of finite crystals allows only two, three, or four. The physics of stable matter allows only three. Two constraints, one answer. The universe is three-dimensional in space, plus one of time from the growth, because those are the only numbers that let a crystal exist and hold matter at the same time.

There is a third hint worth a sentence. When the simulation grows a crystal from scratch and simply measures how many dimensions it acts like, with nothing put in by hand, the answer comes out sharp and stable. The dimension is not imposed. It emerges and it agrees.

\subsection{No hidden dimensions}

Some theories add extra spatial dimensions, curled up too small to see. This model adds none, and that is a feature, not a gap. The classification that allows finite crystals only in two, three, and four dimensions leaves no room for extra spatial dimensions to hide, so there are none to curl up or explain away. Space is three-dimensional, full stop. Time is not a fourth dimension of the same kind, it is the growth of the crystal, a different thing entirely. Where other pictures juggle ten or eleven dimensions, this one has exactly three of space and the one-way growth that is time, and the count is forced rather than chosen.

\textbf{Takeaway:} finite hyperbolic crystals exist only in two, three, or four dimensions, and only three supports stable closed orbits and therefore stable matter. So three spatial dimensions is forced, not assumed.
